%! Change in Arithmetic and Geometric Sequences — Unit 4, Lesson 4.1 — Group Activity (Answer Key)
\documentclass[10pt]{article}
\usepackage{precalculus-article}
\usepackage{precalculus-key}   % pulls in -boxes; do NOT also load -boxes
\newcommand{\TallMath}[1]{$\displaystyle #1\rule[-1.4em]{0pt}{3.2em}$}

\begin{document}

\pageheader{Unit 4, Lesson 4.1}{Group Activity --- Answer Key}
\namepartnerperiod

\vspace{0.08in}

\begin{scenariobox}[Sequences in the Wild]{sky}
\small
Sequences appear in finance, biology, and technology. In this activity you will identify
sequence types, write explicit formulas, and use them to answer real questions. Work through
the tier that matches your current level --- all tiers use the same big ideas.
\end{scenariobox}

\vspace{0.08in}

\begin{tcolorbox}[colback=white, colframe=black!40,
    title=\textbf{Tier R --- Remediate},
    fonttitle=\bfseries, arc=2mm, left=3mm, right=3mm, top=2mm, bottom=2mm]

For each sequence below: (a) identify it as \textbf{arithmetic} or \textbf{geometric},
(b) find $d$ or $r$, (c) write the explicit formula starting at $n = 0$, and
(d) find the 10th term ($n = 10$).

\renewcommand{\arraystretch}{2.0}
\begin{tabularx}{\linewidth}{@{} c >{\centering\arraybackslash}p{3.2cm} >{\centering\arraybackslash}p{2.2cm} >{\centering\arraybackslash}p{2.2cm} >{\centering\arraybackslash}X @{}}
\textbf{Sequence} & \textbf{Type} & $d$ \textbf{or} $r$ & \textbf{Explicit formula} & \textbf{10th term} \\
\midrule
$5,\ 8,\ 11,\ 14,\ \ldots$      & \ans{arithmetic} & \ans{$d=3$} & \ans{$5+3n$}     & \ans{35} \\[6pt]
$3,\ 6,\ 12,\ 24,\ \ldots$      & \ans{geometric}  & \ans{$r=2$} & \ans{$3\cdot2^n$} & \ans{3072} \\[6pt]
$100,\ 90,\ 80,\ 70,\ \ldots$   & \ans{arithmetic} & \ans{$d=-10$} & \ans{$100-10n$}  & \ans{0} \\[6pt]
$1,\ 5,\ 25,\ 125,\ \ldots$     & \ans{geometric}  & \ans{$r=5$} & \ans{$5^n$}       & \ans{$5^{10}=9{,}765{,}625$} \\[6pt]
\end{tabularx}
\end{tcolorbox}

\vspace{0.08in}

\begin{tcolorbox}[colback=white, colframe=black!40,
    title=\textbf{Tier A --- Apply},
    fonttitle=\bfseries, arc=2mm, left=3mm, right=3mm, top=2mm, bottom=2mm]

\begin{enumerate}[itemsep=12pt]

  \item An arithmetic sequence has $a_3 = 12$ and $a_7 = 28$.
        Find the common difference $d$ and write the explicit formula $a_n$.

        \vspace{0.04in}
        $d = \dfrac{a_7 - a_3}{7 - 3} = $ \ans{$\dfrac{28-12}{4} = 4$}

        \vspace{0.04in}
        Explicit formula: $a_n = $ \ans{$a_3 + 4(n-3) = 12 + 4(n-3) = 4n$}

        \vspace{0.04in}
        Check: $a_3 = $ \ans{12} \quad $a_7 = $ \ans{28}

  \item A geometric sequence has $g_2 = 12$ and $g_5 = 96$.
        Find the common ratio $r$ and write the explicit formula $g_n$.

        \vspace{0.04in}
        $r^3 = \dfrac{g_5}{g_2} = \dfrac{96}{12} = $ \ans{8}, so $r = $ \ans{2}

        \vspace{0.04in}
        $g_0 = \dfrac{g_2}{r^2} = $ \ans{$12/4 = 3$}

        \vspace{0.04in}
        Explicit formula: $g_n = $ \ans{$3 \cdot 2^n$}

  \item A phone plan charges \$25 in month 0 and adds \$8 per month.

        \begin{enumerate}[label=(\alph*), itemsep=6pt]
          \item Write the explicit formula for the monthly charge $a_n$ in month $n$:

                $a_n = $ \ans{$25 + 8n$}

          \item Find the charge in month 12: $a_{12} = $ \ans{\$121}

          \item In which month does the charge first exceed \$100? \ans{Month 10 ($a_{10}=\$105$)}
        \end{enumerate}

\end{enumerate}
\end{tcolorbox}

\vspace{0.08in}

\begin{tcolorbox}[colback=white, colframe=black!40,
    title=\textbf{Tier E --- Extend},
    fonttitle=\bfseries, arc=2mm, left=3mm, right=3mm, top=2mm, bottom=2mm]

A bacteria colony triples every hour. After 3 hours there are 810 bacteria.

\begin{enumerate}[itemsep=12pt]

  \item Use the alternate formula $g_n = g_k \cdot r^{(n-k)}$ with $k = 3$, $g_3 = 810$,
        and $r = 3$ to find $g_0$ (the initial population).

        \vspace{0.04in}
        $g_0 = g_3 \cdot r^{(0-3)} = 810 \cdot $ \ans{$3^{-3} = 810/27$} $=$ \ans{30}

  \item Write the explicit formula $g_n$ using $g_0$.

        \vspace{0.04in}
        $g_n = $ \ans{$30 \cdot 3^n$}

  \item Determine the smallest whole-number value of $n$ for which the population first
        exceeds 100{,}000. Show your reasoning.

        \vspace{0.04in}
        We need $g_n > 100{,}000$: \ans{$30 \cdot 3^n > 100{,}000 \Rightarrow 3^n > 3333.3$; $3^7 = 2187 < 3334$, $3^8 = 6561 > 3334$}

        \vspace{0.04in}
        Smallest $n$: \ans{8}

  \item Interpret your answer in context: what does $n$ represent in this problem?

        \vspace{0.04in}
        \ans{$n$ is the number of hours after the colony began. At $n = 8$ hours, the population first exceeds 100{,}000.}

\end{enumerate}
\end{tcolorbox}

\end{document}
